%% =====================================================================
%%  T-03-01  The Five Dangerous Marks
%%  -------------------------------------------------------------------
%%  One idea per file. Core is mandatory. Slots in canonical order.
%%  Max 700 words. No layout commands. No filler. New source material
%%  for this entry goes in sources/<BOOK>/T-03-01.tex -- never by
%%  rewriting this file (docs/RULES.md R0.1).
%% =====================================================================
%% @kind: topic
%% @id: T-03-01
%% @title: The Five Dangerous Marks
%% @subtitle: Typology of the counterpart you should not engage
%% @chapter: c03
%% @order: 10
%% @status: stable
%% @sources: B3
%% @updated: 2026-09-19

\topic{T-03-01}{The Five Dangerous Marks}{Typology of the counterpart you should not engage}

\begin{core}
Five dispositions make a counterpart expensive to work against, and they are recognisable early. The instruction is not to handle them better but to decline the engagement: whatever you hope to get from them is unlikely to be worth what they will cost.
\end{core}
\begin{mechanism}
Each type defeats a different assumption. The ordinary model of a counterpart --- someone who responds to incentives, forgets small injuries, and reacts in proportion --- fails in a specific way for each. Pride reacts out of all proportion and immediately. Fragile self-esteem reacts slowly and indefinitely. Suspicion misreads everything, including your honesty, and is therefore easy to fool and dangerous to be suspected by. Cold calculation waits and remembers. Literal inability to follow a scheme wastes your resources rather than defeating them. Since the failure modes are structural rather than situational, they cannot be managed by better technique.
\end{mechanism}
\begin{tells}
  \item \emph{Pride}: reacts violently to a slight that was minor or unintended, and the reaction cannot be reasoned with.
  \item \emph{Fragile self-esteem}: no open reaction, then a long series of small obstructions you only notice in aggregate.
  \item \emph{Suspicion}: assumes the worst of everyone, reads motives into accidents, and is easily led to distrust a third party.
  \item \emph{Long memory}: shows nothing at the time, is cold and unaffectionate, and settles accounts later from a position of advantage.
  \item \emph{Literal}: does not take a hint, does not follow an implication, and reacts to a joke or a story with flat literalness.
\end{tells}
\begin{moves}
  \item Test early and cheaply, before committing anything --- a story, a joke, a small favour, a small slight.
  \item Read the reaction against the five patterns above rather than against what you expected.
  \item If a pattern matches, withdraw while the cost of withdrawing is still small.
  \item If you have already injured one of them, either settle it completely or remove yourself from their reach. Partial settlement is worse than none.
\end{moves}
\begin{counters}
  \item Assume you may be the mark in someone else's typology, and ask which pattern you fit.
  \item Notice disproportionate reactions in yourself. They are the same tells.
  \item Where you cannot decline the engagement --- a relative, a manager --- reduce the surface area: fewer commitments, fewer confidences, nothing in writing that can be held.
  \item Do not attempt to deceive the literal type; the cost is wasted time, and the return is nothing.
\end{counters}
\begin{cost}
The typology is a heuristic built on extremes, and applied loosely it becomes a way of writing off anyone who has ever been difficult. Most people are none of these types most of the time; the categories describe a settled disposition, not a bad day. It also carries a survivorship problem --- the types are named by the people who beat them, so the quiet counterpart who simply declined to engage never appears in the list at all.
\end{cost}

\sources{B3}
\seesources
\seealso{T-03-02, T-08-01, T-08-02}
